\scanpage{6}

\deleted{\uline{Corollaire} 4.1. ter}\inserted{Proposition A} Soit $f:X\quad Y$ un morphisme de présentation finie, avec $Y$ \deleted{\unsure{réduit}}\inserted{$\neq\varnothing$}, $F$ un Module de présentation finie sur $X$. Alors, quitte à remplacer $Y$ par un ouvert non vide convenable, on a ce qui suit: il existe un sous-schéma fermé \deleted{\unsure{X de}} $Z$ de $Y$, \inserted{de présentation finie sur $Y$} tel que pour tout changement de base $Y'\inserted{\to}Y$, la condition nécessaire et suffisante pour que $F'$ soit plat par rapport à $Y'$ est que $Y'\quad Y$ se factorise \deleted{par} en $Y'\quad Z\quad Y$.
De plus,\inserted{si $Y$ est affine et $X$ est réunion finie d'ouverts affines $U_i$, alors}(quitte à remplacer $Y$ par un ouvert non vide plus petit), on a ceci: si la condition précédente est vérifiée pour $Y'$, et si $Y'$ est affine, \inserted{d'anneau $A'$} alors \deleted{pour tout ouvert affine $U$ de $X$,}\inserted{pour tout $i$,} $\Gamma(U_i,F')$ est un $A'$-module \uline{libre}.

On peut évidemment supposer $Y$ affine, i.e. $Y=\operatorname{Spec}(A)$. Notant que la situation $(X,F)$ sur $Y$ provient par changement de base d'une situation analogue $(X_o,F_o)$ sur un $Y_o=\operatorname{Spec}(A_o)$, où $A_o$ est un sous-anneau noethérien de $A$, on est ramené trivialement au cas où $A$ est noethérien. Soit $I$ l'idéal des éléments nilpotents de $A$, $A_o=A/I$, $Y_o=\operatorname{Spec}(A_o)$, et dénotons par un ${}_o$ l'effet du changement de base $Y_o\quad Y$. En vertu du théorème de platitude générique (par 9), quitte à remplacer $Y$ par un ouvert non vide convenable, on peut supposer que les $M_{io}=\Gamma(U_{i\,o},F_{i\,o})=M_i\otimes_A A_o$ (où $M_i=\Gamma(U_i,F)$) sont des $A_o$-modules \uline{libres}. A partir de là, l'assertion résulte facilement de 4.1. bis : pour tout $i$, on trouve un idéal $J_i$ dans $A$, tel que pour \unsure{toute} $A$-algèbre $A'$, $M_i\otimes_A A'$ soit libre, ou encore plat, sur $A'$, sss $A\quad A'$ est nul sur $J_i$. Soit $J$ l'idéal somme des $J_i$, alors pour tout $i$, $M_i\otimes_A A'$ est libre sur $A'$, ou encore plat sur $A'$, sss $A\quad A'$ s'annule sur $J$, i.e. $A\quad A'$ se factorise en $A\quad A/J\quad A'$. Il suffit maintenant de prendre $Z=V(J)=\operatorname{Spec}(A/J)$, notant que \deleted{\unsure{l'on}} pour vérifier que $Z$ satisfait à la condition de l'énoncé, il suffit de le vérifier pour des $Y'$ affines. \deleted{Noter d'ailleurs}

\mspage{5}
